\chapter{Hạn chế và hướng phát triển}
\label{chap:limits}

Nhóm liệt kê các hạn chế một cách trung thực, để người đọc phân biệt phần đã chạy thật với phần còn ở mức demo hoặc thiết kế.

\section{Hạn chế hiện tại}

\begin{enumerate}
  \item \textbf{BR1--BR3 vẫn chạy trong runtime demo.} Rule engine thông báo, chính sách chấm bài và luật gợi ý đã hoàn chỉnh và có test, nhưng các dịch vụ tiêu thụ Kafka và adapter Redis cho môi trường production vẫn là stub. Sandbox demo là tiến trình con và SQLite cục bộ, không phải ranh giới bảo mật.
  \item \textbf{BR2 chưa có nhánh đặc trưng phiên và mô hình học máy.} Đề bài yêu cầu kết hợp đặc trưng lô và đặc trưng phiên qua feature store. Hiện mới có bảng đặc trưng lô trong lakehouse; việc đẩy bảng này lên Redis và nhánh streaming chưa được hiện thực.
  \item \textbf{Độ trễ streaming ở mức giây.} Spark Structured Streaming cho cảnh báo BR5 sau khoảng 10 giây và lớp phủ sau khoảng 7 giây, cao hơn mức dưới một giây mà Flink có thể đạt. Đây là đánh đổi có chủ đích (ADR 0009, 0014).
  \item \textbf{Embedding mang tính từ vựng.} Bộ nhúng băm không hiểu ngữ nghĩa (``join'' không khớp ``merge''). Đã có cổng \code{Embedder} và đường reindex để thay bằng mô hình thật.
  \item \textbf{Lớp phủ không mang phiên bản nội dung.} Vào ngày một video được quay lại, lớp phủ trộn hai phiên bản; vì vậy nó chỉ được hiển thị cho phiên bản mới nhất.
  \item \textbf{Lịch sử BR5 trong kho dữ liệu đến từ dữ liệu tổng hợp.} Runtime BR4/BR5 lưu điểm và cảnh báo trong PostgreSQL; cần phát chúng lên \code{risk.learner_score.v1} và \code{risk.dropout_alert.v1} (hoặc thêm bảng vào connector CDC) để kho dữ liệu nhận từ đường thật.
  \item \textbf{Hiệu năng BR4.} Mỗi lần sửa dữ liệu làm mới toàn bộ bản tổng hợp bị ảnh hưởng và dùng phép nối JSONB; với 100.000 lượt chấm, một lần sửa tốn khoảng 4 giây.
  \item \textbf{Vận hành đơn nút, bảo mật cục bộ.} Trino, OpenSearch, Kafka chạy một nút; plugin bảo mật của OpenSearch tắt; token tĩnh thay cho hệ thống định danh của trường. Chưa có giám sát (độ trễ consumer, thời gian checkpoint, SLA độ tươi).
  \item \textbf{Dữ liệu tổng hợp, không nhãn thật.} Không có nhãn bỏ học thật nên không đánh giá được độ chính xác của cảnh báo; các số đo hiệu năng là quan sát trên một máy.
  \item \textbf{Lệch dữ liệu mẫu.} Seed của lakehouse khai báo khoảng 18 khoá học và video không có trong cơ sở dữ liệu LMS mẫu, nên job reindex báo các mục này là lỗi.
\end{enumerate}

\section{Hướng phát triển}

\begin{itemize}
  \item Hoàn thiện các dịch vụ Kafka cho BR1 và BR3, và adapter Redis cho dedup và giới hạn tần suất.
  \item Bổ sung nhánh đặc trưng phiên cho BR2 (cửa sổ phiên 30 phút), đẩy \code{features.learner_longterm} lên Redis, và thử một mô hình gợi ý học từ dữ liệu, so sánh với luật hiện tại.
  \item Chuyển các job cần độ trễ dưới một giây sang Flink khi nền tảng lên Flink 2.x và connector Delta Kernel ổn định.
  \item Thay bộ nhúng băm bằng mô hình sentence-embedding, chạy một lần reindex bắt buộc.
  \item Chuyển các bảng snapshot (\code{fct_enrollment}, \code{fct_lesson_progress}) sang MERGE dựa trên Change Data Feed khi khối lượng đủ lớn; với dashboard dưới giây ở quy mô lớn, gắn StarRocks vào cùng Hive Metastore.
  \item Triển khai lên Kubernetes với nhiều nút, TLS và xác thực tập trung; giám sát bằng Prometheus cùng Grafana; bổ sung quy trình xoay khoá bút danh.
  \item Thu thập nhãn kết quả thật để hiệu chỉnh và đánh giá điểm rủi ro BR5.
\end{itemize}
